\subsection{A step over a truncated parent}\label{cc:uf:sec:parent-step}
Throughout this section the parent is the truncated velocity to which
we add the new wave. We prescribe the wave vector and seed at the origin
and extend the seed to the labels in the envelope.
Section~\ref{cc:sec:iteration} verifies the two assumptions
\eqref{cc:uf:tower:eq:wave-origin} on the central pair for the first
wave and for each later wave.

We use the Fourier modes $e^{2\pi i\xi\cdot x}$ for $\xi\in\mathbb Z^3$.
A field is \emph{supported on frequencies at most $\Lambda$} if its
Fourier coefficients vanish for $|\xi|>\Lambda$. Let
the parent $U$ be a smooth odd divergence-free field on
$[0,T]\times\mathbb T^3$, $T\le1$, and use its flow $X$, the matrix $M$
and the operators $\mathbf d$, $\widetilde{\mathbf d}$ of
Subsection~\ref{cc:sec:material}. Let $\mathsf D\ge2$ be a number such
that $U$ is supported on frequencies at most $2\mathsf D$ and
\begin{equation}\label{cc:uf:tower:eq:parent-bound}
 \norm{\nabla U}_\infty\le\mathsf D,\qquad
 \sup|(D_aX)^{\pm1}|\le\mathsf D .
\end{equation}
Since $U$ is odd, the flow fixes the origin, and $D_aX$ and $M$ are even
in $a$. Let $p_U$ be the function of zero mean with
$\Delta p_U=-\tr((\nabla U)^2)$, and let
$R_U=U_t+U\cdot\nabla U+\nabla p_U$ be the residual of the parent; no
other equation is assumed. Thus $R_U=U_t+\PP\diver(U\otimes U)$ is odd,
divergence free and supported on frequencies at most $4\mathsf D$.
To apply the endpoint inverse of Theorem~\ref{cc:uf:mean:inverse}
to the mean correction over $U$, assume that
\begin{equation}\label{cc:uf:tower:eq:parent-coercive}
 D^2p_U-\nabla R_U\le\Id\qquad\hbox{on }[0,T]\times\mathbb T^3
\end{equation}
as quadratic forms. By \eqref{cc:uf:mean:eq:Hsharp} this is the
hypothesis $H^\#\le\Id$ of that theorem.

The wave is that of Subsection~\ref{cc:sec:material} over $U$, with its
constant vector $k_*\ne0$ and transported field
$\kappa=(D_aX)^{-T}k_*$, its activation time $t_0\in[0,T)$ and its cutoff
in time $\widetilde\chi$. Let $\lambda\ge1$ be the frequency of the wave:
its wave vector at the origin has length $|\kappa(t_0,0)|=\lambda$ at the
activation time. Since $k_*=(D_aX)^T\kappa$ at every time and label, the
bound \eqref{cc:uf:tower:eq:parent-bound} for the deformation gives
\begin{equation}\label{cc:uf:tower:eq:wave-vector}
 \mathsf D^{-1}\lambda\le|k_*|\le\mathsf D\lambda,\qquad
 \mathsf D^{-2}\lambda\le|\kappa|\le\mathsf D^2\lambda .
\end{equation}
The profile is $F_0=F_h$ of Lemma~\ref{cc:uf:exact:lem:profile}, with
$h=\lambda^{-1/2}$. At the origin the seed $\mathbf a(t,0)$,
$t_0\le t\le T$, is a solution of the seed equation
\eqref{cc:uf:step:eq:seed} at the label $a=0$, perpendicular to
$\kappa(t,0)$ and normalized by
$\max_{t_0\le t\le T}|\mathbf a(t,0)|=\lambda^{-3/4}$. It does not
vanish, because it is a nonzero solution of a linear ordinary
differential equation. Put
\[
 \gamma(t)=\lambda^{3/4}|\mathbf a(t,0)|,\qquad
 \nu(t,0)=\Bigl(\frac{\kappa}{|\kappa|}\cdot
          M\frac{\mathbf a}{|\mathbf a|}\Bigr)(t,0)
 \qquad(t_0\le t\le T).
\]
As in Lemma~\ref{cc:lem:frame}, $\nu(t,0)$ is the first coupling of the
pair $(\kappa,\mathbf a)$ at the origin over the parent, and
$|\mathbf a(t,0)||\kappa(t,0)|$ is its strength. At the other labels
the seed is the solution of \eqref{cc:uf:step:eq:seed} on $[t_0,T]$ with
\begin{equation}\label{cc:uf:tower:eq:start}
 \mathbf a(t_0,a)=\bigl(\Id-P_{\kappa(t_0,a)}\bigr)\mathbf a(t_0,0),
 \qquad P_\kappa=\frac{\kappa\otimes\kappa}{|\kappa|^2},
\end{equation}
the projection of the vector at the origin onto the plane perpendicular
to $\kappa(t_0,a)$. It is even in $a$.

Fix a smooth function $\chi$ on $[0,\infty)$ with $0\le\chi\le1$,
$\chi=1$ on $[0,1]$ and $\chi=0$ on $[2,\infty)$, and take the even
envelope $g(a)=\chi(|a|/\ell_*)$, equal to one on the ball of radius
$\ell_*$ specified below. The leading oscillation is the field
$A_1=g\mathbf aF_h$ of \eqref{cc:uf:hier:eq:primary}, and it enters
through $V_1=\widetilde\chi A_1$ of \eqref{cc:uf:hier:eq:switched}.

At the origin, assume that for $t_0\le t'\le t\le T$,
\begin{equation}\label{cc:uf:tower:eq:wave-origin}
 \frac{|\mathbf b(t)|}{\gamma(t)}
 \le\mathsf D\,\frac{|\mathbf b(t')|}{\gamma(t')},\qquad
 \nu(t,0)\ge\mathsf D^{-1},
\end{equation}
the first for every solution $\mathbf b$ of the seed equation at the
origin that is perpendicular to $\kappa(\cdot,0)$. The first inequality
bounds the growth of every transverse solution by that of the seed,
up to the factor $\mathsf D$. On a sufficiently small envelope this
gives the weighted estimate for the transverse corrections, and the
second inequality keeps $\kappa\cdot M\mathbf a$ positive there
(Lemma~\ref{cc:uf:tower:lem:inverse}). With this sign, the lower bound
$F_h'\ge-Ch$ controls the leading pressure tensor in
\eqref{cc:uf:exact:eq:primary-hess} from above.
Both inequalities can be imposed before choosing the frequency and the
depth of the expansion.

On the envelope, \eqref{cc:uf:tower:eq:near-origin} bounds the seed
by twice its size at the origin. Since $\|F_h\|_\infty\le Ch$, the leading
velocity has size at most $C\lambda^{-3/4}h=C\lambda^{-5/4}$.
The first mean equation removes both the quadratic wave forcing and
$R_U$, so we require the parent residual to be at most the square of
this scale,
\begin{equation}\label{cc:uf:tower:eq:parent-small}
 \norm{R_U}_{C_tL^\infty}\le\lambda^{-5/2}.
\end{equation}
The radius of the envelope is a negative power of $\mathsf D$ with an
absolute exponent, $\ell_*\le\mathsf D^{-3}$. The exponent is
fixed sufficiently large for the finitely many estimates of
Subsections~\ref{cc:uf:trunc:sec:parent}--\ref{cc:uf:sec:transverse}.
In particular the support of $g$ lies in the interior of
$[-1/2,1/2]^3$.

Constants and exponents written $C$ are absolute: they do not depend on
the parent, the wave, an order of differentiation or the depth of the
expansion. A constant written $C_n$ depends only on the order $n$, on
the fixed profiles, and on the constants of the same kind in the
hypotheses of the statement in which it appears; $C_{n,p}$ and $C_{m,P}$
depend also on a level $p$ or on the depth $P$. These constants may
change from one estimate to the next, and no bound on their growth in
$n$ is used. In the iteration these constants and the powers of
$\mathsf D$ are at most a small power of $\lambda$; see
\eqref{cc:par:step}.
